Este sistema de monitoramento tem como meta desenvolver uma aplicação para auxiliar os produtores criarem camarões de boa qualidade e manter o controle do cativeiro. Dentre os objetivos estão:

\begin{enumerate}

\item Monitorar remotamente os cativeiros, acompanhando pH, temperatura e amônia da água para garantir a qualidade do ambiente e alertar o usuário sobre desvios nos parâmetros.

\end{enumerate}

\section*{Objetivos específicos}

\begin{enumerate}
    \item Registrar e armazenar dados históricos dos cativeiros para possibilitar análises comparativas e avaliação da qualidade ao longo do tempo.

    \item Automatizar o fornecimento de ração, controlando horários e quantidades conforme parâmetros definidos pelo usuário.

    \item Fornecer relatórios e gráficos gerados a partir das informações monitoradas, permitindo uma análise detalhada e fundamentada das condições do cultivo.
    
    \item Cadastrar usuários, cativeiros e pontos de monitoramento, facilitando a gestão e o acesso centralizado às informações.
    
    \item Desenvolver uma aplicação móvel intuitiva e acessível, proporcionando uma experiência prática e eficiente ao carcinicultor.

\end{enumerate}